% Level 2, Week 6 -- Review and Checkpoint (student workbook)
\chapterweek{6}{Review and Checkpoint}{Bugs are normal: detectives find them with a strategy, not by luck.}

\section{What you will learn}
\begin{itemize}[leftmargin=1.2em]
  \item How to read a \textbf{traceback} from the bottom up.
  \item Three more detective tools: \textbf{print-debugging}, the \textbf{rubber duck} and the \textbf{trace table}.
  \item How to write a short \textbf{bug report}, so you (and others) remember what you learned.
\end{itemize}

\section{The big idea: debug like a detective}
Every programmer writes bugs. Good programmers find them with a \textbf{strategy}: they collect clues, change \textbf{one} thing and test again.

\subsection{Tool 1: read the traceback from the bottom up}
When a program crashes, Thonny shows a red \textbf{traceback} like this one:
\begin{outputcode}
Traceback (most recent call last):
  File "starter.py", line 24, in <module>
  File "starter.py", line 17, in make_card
  File "starter.py", line 10, in card_line
    return label + ": " + value
TypeError: can only concatenate str (not "int") to str
\end{outputcode}
\begin{enumerate}[leftmargin=1.6em]
  \item \textbf{WHAT} went wrong: the \textbf{last line}. Here: a \texttt{TypeError}, text + number.
  \item \textbf{WHERE}: the \texttt{File} line just above it. Here: line 10, in the function \texttt{card\_line}.
  \item \textbf{HOW} it got there: the lines further up. The main program (line 24) called \texttt{make\_card}, which called \texttt{card\_line}.
\end{enumerate}
In Thonny you can click the blue link to jump straight to that line.

\begin{center}
\renewcommand{\arraystretch}{1.3}
\begin{tabularx}{\linewidth}{@{}lX@{}}
\toprule
\textbf{Error name} & \textbf{It usually means\ldots} \\ \midrule
\texttt{TypeError} & two types that do not go together, like text + number \\
\texttt{IndexError} & a list index that does not exist (remember: the last index is \texttt{len - 1}) \\
\texttt{KeyError} & a dictionary key that does not exist (check the spelling) \\
\texttt{NameError} & a variable name Python does not know (typo, or missing quotes) \\
\texttt{ValueError} & the right type but a bad value, like \texttt{int("banana")} \\ \bottomrule
\end{tabularx}
\end{center}

\subsection{Tool 2: print-debugging}
Not sure what is inside a variable? Ask Python! Add a \texttt{print()} line, run, look, then remove it.
\begin{pycode}
print("best is", best, "length", len(best))
\end{pycode}

\subsection{Tool 3: the rubber duck}
Explain your code, line by line, out loud, to a partner, a toy or a rubber duck. Say what each line \emph{should} do.
Very often you hear yourself say the bug before you reach the end. This is the best tool for bugs that do not crash.

\subsection{Tool 4: the trace table}
Follow the code with a pencil, like a very slow computer. Write a column for each variable and a row each time round the loop.
When the table shows a value you did not expect, you have found the bug.

\begin{tip}[Bugs are normal]
If you are stuck, you are not bad at coding: you are in the middle of learning. Say \enquote{I can't fix this \emph{yet}}, pick a tool and try one thing.
\end{tip}

\section{Bug-hunt workshop}
Each station has a folder in \texttt{exercises/week06}. Open \texttt{starter.py}, find the bug, fix it, then check with \texttt{check\_my\_work.py}.
Fill in a bug report for each bug you fix.

\begin{exercise}[Station 1: Hero Card (TypeError)]
Folder: \texttt{ex-bug-hero-card}. We do this one together.\par
Error name: \rule{4cm}{0.4pt}\quad Line: \rule{1.5cm}{0.4pt}\quad Tool I used: \rule{3.5cm}{0.4pt}\par
I expected: \rule{5cm}{0.4pt}\quad It really did: \rule{5cm}{0.4pt}\par
What I changed: \answerlines{1}
\end{exercise}

\begin{exercise}[Station 2: High Scores (IndexError, 2 bugs)]
Folder: \texttt{ex-bug-high-scores}. Draw the list with its index numbers. Try \texttt{print("best is", best)} inside \texttt{podium()}.\par
Bug 1 -- error name: \rule{3.5cm}{0.4pt}\quad line: \rule{1.5cm}{0.4pt}\quad what I changed: \rule{5cm}{0.4pt}\par
Bug 2 -- error name: \rule{3.5cm}{0.4pt}\quad line: \rule{1.5cm}{0.4pt}\quad what I changed: \rule{5cm}{0.4pt}
\end{exercise}

\begin{exercise}[Station 3: Rocket Countdown (infinite loop)]
Folder: \texttt{ex-bug-countdown}. The safety net stops the stuck loop. Fill in the trace table for \texttt{countdown(5)} as the starter is now:
\begin{center}
\renewcommand{\arraystretch}{1.4}
\begin{tabular}{|c|c|p{4cm}|}
\hline \textbf{round} & \textbf{n} & \textbf{numbers} \\ \hline
1 & 5 & [5] \\ \hline
2 & & \\ \hline
3 & & \\ \hline
\end{tabular}
\end{center}
What should change each time round, but does not? \rule{6cm}{0.4pt}\par
What I changed: \answerlines{1}
\end{exercise}

\begin{exercise}[Station 4: Loot Shop (3 sneaky bugs, no crash)]
Folder: \texttt{ex-bug-loot-shop}. First work it out on paper for the basket \texttt{[12, 5, 30]}:\par
Total: \rule{1.5cm}{0.4pt}\quad Cheapest: \rule{1.5cm}{0.4pt}\quad Can you buy it all with 50 coins? \rule{1.5cm}{0.4pt}\par
Now run it and explain each function to your duck. The three bugs I fixed:
\answerlines{3}
\end{exercise}

\begin{exercise}[Station 5: Paper detectives (no computer)]
This function should give back the biggest number in a list. Fill in the trace table for \texttt{biggest([-5, -2, -8])}.
\begin{center}
\begin{minipage}{0.45\linewidth}
\begin{pycode}
def biggest(numbers):
    best = 0
    for n in numbers:
        if n > best:
            best = n
    return best
\end{pycode}
\end{minipage}\hfill
\begin{minipage}{0.45\linewidth}
\renewcommand{\arraystretch}{1.4}
\begin{tabular}{|c|c|c|}
\hline \textbf{n} & \textbf{n > best?} & \textbf{best} \\ \hline
start & -- & 0 \\ \hline
-5 & & \\ \hline
-2 & & \\ \hline
-8 & & \\ \hline
\end{tabular}
\end{minipage}
\end{center}
It gives back \rule{1.5cm}{0.4pt}, but the biggest number is \rule{1.5cm}{0.4pt}. What is the bug, and how would you fix it?
\answerlines{2}
\end{exercise}

\begin{vocab}
  \vocabword{bug}{A mistake in a program that makes it crash or give the wrong answer.}
  \vocabword{debugging}{Finding and fixing bugs, step by step, with a strategy.}
  \vocabword{traceback}{The red error message. Read it from the bottom up: what, then where.}
  \vocabword{print-debugging}{Adding \texttt{print()} lines to see what is really inside your variables.}
  \vocabword{rubber duck debugging}{Explaining your code line by line out loud, to find the bug yourself.}
  \vocabword{trace table}{A table where you follow code by hand and write each variable's value.}
  \vocabword{bug report}{A short note: what went wrong, where, and what you changed.}
\end{vocab}

\begin{reflection}
Six weeks ago I could not \ldots{} and now I can \ldots
\answerlines{2}
Which detective tool will you try first next time? Why?
\answerlines{2}
\end{reflection}
